\documentclass[polish]{article}
\usepackage{graphicx} % Required for inserting images

\usepackage[polish]{babel}
\usepackage{polski}
\usepackage{placeins}
\usepackage{dirtytalk}
% Set page size and margins
% Replace `letterpaper' with `a4paper' for UK/EU standard size
\usepackage{csquotes}
% Useful packages
\usepackage{float}
\usepackage{amsmath}
\usepackage{graphicx}
\usepackage{hyperref}
\usepackage{graphicx}
\usepackage{listings}
\usepackage{xcolor}
\usepackage{amssymb}
\usepackage{float}
\definecolor{codegreen}{rgb}{0,0.6,0}
\definecolor{codegray}{rgb}{0.5,0.5,0.5}
\definecolor{codepurple}{rgb}{0.58,0,0.82}
\definecolor{backcolour}{rgb}{0.95,0.95,0.92}

\lstdefinestyle{mystyle}{
    backgroundcolor=\color{backcolour},   
    commentstyle=\color{codegreen},
    keywordstyle=\color{magenta},
    numberstyle=\tiny\color{codegray},
    stringstyle=\color{codepurple},
    basicstyle=\ttfamily\footnotesize,
    breakatwhitespace=false,         
    breaklines=true,                 
    captionpos=b,                    
    keepspaces=true,                 
    numbers=left,                    
    numbersep=5pt,                  
    showspaces=false,                
    showstringspaces=false,
    showtabs=false,                  
    tabsize=2
}
\lstset{style=mystyle}
\hypersetup{
    colorlinks=true,
    linktoc=all     %set to all if you want both sections and subsections linked
}
\usepackage{mdframed}
\title{Zagadnienia do 2. kolokwium}
\author{}
\date{}

\begin{document}

\maketitle

\section{Zagadnienia}
\begin{enumerate}
    \item Funkcje tworzące momenty.
    \item Rozkłady omawiane na wykładzie, m.in. dwumianowy, geometryczny, Poissona, jednostajny, wykładniczy, normalny, gamma, $\chi^2$, $t$ Studenta, $F$ i ich własności.
    \item Próba losowa, statystyki, średnia i wariancja w prostej próbie losowej, rozkład średniej, estymatory.
    \item Naiwny klasyfikator bayesowski.
    \item Zmienne losowe wielowymiarowe (rozkład łączny, rozkłady brzegowe, rozkłady warunkowe, niezależność, kowariancja, współczynnik korelacji, macierz kowariancji, momenty i ich własności). 
    \item Prawa wielkich liczb, Centralne Twierdzenie Graniczne, 
    \item Metoda największej wiarygodności, metoda momentów.
    \item Konstrukcja przedziałów ufności.
    \item Testy statystyczne. W przypadku testów warto spojrzeć na język R:  $t.test$, $chisq.test$, $var.test$. 
    \item Regresja liniowa. 
\end{enumerate}

\section{Dodatkowe materiały}
\begin{itemize}
    \item \href{https://rpsychologist.com/correlation/}{Współczynnik korelacji}
    \item \href{https://rpsychologist.com/likelihood/}{Metoda największej wiarygodności}
    \item \href{https://rpsychologist.com/d3/tdist/}{Rozkład $t$}
    \item \href{https://m.youtube.com/watch?v=zeJD6dqJ5lo}{Centralne Twierdzenie Graniczne}
    \item \href{https://m.youtube.com/watch?v=vemZtEM63GY}{$p$-wartości}
    \item \href{https://m.youtube.com/watch?v=VekJxtk4BYM}{$t$-test}
    \item \href{https://m.youtube.com/watch?v=MntX3zWNWec}{Prawo wielkich liczb}
\end{itemize}

\section{Opis zadań}
Zadania będą zbliżone do zadań egzaminacyjnych. Na pewno pojawi się jedno zadanie dotyczące wielowymiarowych zmiennych losowych (podobne do zadań dostępnych na platformie UPEL). Jest bardzo duża szansa na zadanie związane z estymacją, przedziałami ufności lub testami statystycznymi. Prawdopodobnie będą cztery zadania. Jedno zadanie będzie mniej standardowe.

\section{Przykładowe zadania}
\begin{enumerate}
    \item Opisz konstrukcję przedziału ufności dla różnicy wartości średnich w dwu równej długości próbach z rozkładów normalnych o nieznanych wariancjach.
    \item Dla testu $t$ dla wartości średniej w pojedynczej próbie z rozkładu normalnego o nieznanej wariancji z dwustronną hipotezą alternatywną odpowiedz na następujące pytania.

\begin{itemize}
    \item[a)] Jaka jest postać statystyki testowej?
    \item[b)] Jaki rozkład ma ta statystyka jeśli spełniona jest $H_0$?
    \item[c)] Jak się konstruuje zbiór krytyczny na poziomie istotności $\alpha$?
    \item[d)] Jak się wylicza $p$-wartość jako funkcję obliczonej wartości statystyki testowej?
\end{itemize}
    \item Na następnej stronie znajduje się raport z testu statystycznego. Dane pochodzą z rozkładu normalnego. Hipoteza zerowa mówi, że rozkład próby wartość średnią 1, rozważamy dwustronną hipotezę alternatywną.

\begin{itemize}
    \item[a)] Podaj definicję statystyki testowej używanej w tym teście. Jaki rozkład prawdopodobieństwa ma ta statystyka?
    \item[b)] Podaj sposób konstrukcji przedziału ufności wymienionego w raporcie.
    \item[c)] Podaj sposób obliczenia $p$-wartości wymienionej w raporcie.
    \item[d)] Zinterpretuj wyniki testu.
\end{itemize}

\begin{verbatim}
##
## One Sample t-test
##
## data:  x
## t = -1.908, df = 99, p-value = 0.05929
## alternative hypothesis: true mean is not equal to 1
## 95 percent confidence interval:
##  -0.8244266  1.0357342
## sample estimates:
## mean of x
## 0.1056538
\end{verbatim}
\item Wyznacz funkcję generującą momenty rozkładu geometrycznego. Przy jej pomocy oblicz wartość oczekiwaną i wariancję tego rozkładu.
\item Niech zmienna losowa dwuwymiarowa $(X, Y)$ ma rozkład jednostajny na trójkącie $K$ o wierzchołkach $(0, 0)$, $(2, 0)$, $(2, 1)$, tzn. rozkład ciągły o gęstości
\[
f(x, y) =
\begin{cases} 
c & \text{dla } (x, y) \in K, \\
0 & \text{dla } (x, y) \notin K,
\end{cases}
\]
dla pewnej stałej $c \in \mathbb{R}$. Wyznacz wartość $c$. Wyznacz wartość oczekiwaną i macierz kowariancji $(X, Y)$. Wyznacz współczynnik korelacji $\varrho_{X,Y}$. Wyznacz rozkłady brzegowe względem obu składowych, gęstość warunkową $f(y|x)$ i prawdopodobieństwo
\[
P\left( |Y| < \frac{1}{2} \ \middle| \ X = \frac{3}{2} \right).
\]
Sprawdź niezależność zmiennych $X$ i $Y$.

\item Podaj wypowiedź twierdzenia Lindeberga-Lévy'ego. Jaka jest definicja typu zbieżności występującej w tym twierdzeniu? Podaj podstawowe zastosowania tego twierdzenia w statystyce.

\end{enumerate}

\section{Mniej standardowe zadania (z egzaminów)}
\begin{enumerate}
\item Udowodnij, że jeśli zmienne losowe $X_1$ i $X_2$ są niezależne, przy czym $X_1 \sim Exp\left(\frac{1}{2}\right)$ i $X_2 \sim Exp\left(\frac{1}{2}\right)$, to $X_1 + X_2 \sim \chi^2(4)$.

\item Udowodnij, że jeśli wariancja zmiennej losowej jest równa $0$, to zmienna ta jest stała z prawdopodobieństwem $1$.
\item Udowodnij, że dla dowolnej zmiennej dwuwymiarowej \((X, Y)\) i dla dowolnych stałych \(a, b, c, d \in \mathbb{R}\) zachodzi równość
\[
\varrho_{aX+b, cY+d} = \mathrm{sgn}(ac) \varrho_{X, Y},
\]
gdzie \(\varrho_{U, V}\) oznacza współczynnik korelacji \(U\) i \(V\).
\item Udowodnij, że dla dowolnej zmiennej dwuwymiarowej \((X, Y)\) zachodzi równość 
\[
\text{Cov}(X, Y) = \mathbb{E}(XY) - \mathbb{E}(X)\mathbb{E}(Y).
\] 

\end{enumerate}
\end{document}

